% Compilador recomendado no Overleaf: XeLaTeX. Também compatível com pdfLaTeX.
\documentclass[12pt,a4paper]{article}
\usepackage{iftex}
\ifPDFTeX\usepackage[utf8]{inputenc}\fi
\usepackage[T1]{fontenc}
\usepackage[brazil]{babel}
\usepackage{mathptmx}
\usepackage[left=2cm,right=3cm,top=3cm,bottom=2cm,headheight=15pt]{geometry}
\usepackage{graphicx,booktabs,array,longtable,setspace,needspace,fancyhdr}
\usepackage{xurl}
\usepackage[hidelinks,unicode]{hyperref}
\urlstyle{same}
\setlength{\parindent}{1.25cm}
\setlength{\parskip}{3pt}
\setlength{\LTpre}{2pt}\setlength{\LTpost}{4pt}
\setlength{\emergencystretch}{2em}
\onehalfspacing
\pagestyle{fancy}\fancyhf{}\fancyhead[R]{\fontsize{10}{12}\selectfont\thepage}
\renewcommand{\headrulewidth}{0pt}
\newcommand{\heading}[1]{\par\Needspace{4\baselineskip}\noindent\textbf{#1}\par\vspace{5pt}}
\newcommand{\captiontext}[1]{\par\begingroup\fontsize{10}{12}\selectfont%
\setstretch{1}\noindent #1\par\endgroup\vspace{3pt}}
\hypersetup{pdfauthor={Joey Alan de Freitas Solis; Paulo Henrico}}
\begin{document}

\begin{center}\large\textbf{Acompanhamento e auditoria final - ENEM 2023}\end{center}
Joey Alan de Freitas Solis e Paulo Henrico


Universidade de Fortaleza (Unifor) - Professor Tulio


Revisão de 7 de novembro de 2026. Este documento registra desenvolvimento, execuções conferidas e limites da validação. Não comprova apresentação oral, submissão no AVA ou compilação dentro de uma conta do Overleaf.


\heading{1 Objetivo e metodologia}
Estimar NU\_NOTA\_CN a partir de quarenta atributos de perfil socioeconômico, escolar e demográfico, comparando sete famílias de regressão. A fonte nacional contém 3.933.955 linhas e 76 colunas. Foram mantidos os 2.692.427 alvos observados, incluindo 16.547 notas zero. As 1.241.528 ausências de alvo não foram imputadas.


A divisão 80/20, semente 42, reserva 538.486 participantes para teste. Os 2.153.941 restantes usam três dobras de validação cruzada, iguais para os modelos. Transformações são ajustadas no treino de cada dobra; identificadores, presença, notas e resultados de outras avaliações ficam fora dos preditores. O teste foi aberto somente após seleção e congelamento das receitas.


Os dois ZIPs entregues pelo usuário estão em \nolinkurl{dados/raw/}. Os arquivos passaram por CRC; o ZIP nacional e o CSV coincidem com os hashes históricos. A comparação de todas as células normalizadas do CSV com os 79 Parquets existentes foi aprovada. O complemento de competências da redação está preservado, sem integração à matriz de atributos.


\heading{2 Checkpoints e entregáveis}
\begingroup\singlespacing\setlength{\tabcolsep}{3pt}
\renewcommand{\arraystretch}{1.08}
\begin{longtable}{@{}>{\raggedright\arraybackslash}p{\dimexpr.46\linewidth-3.0pt\relax}>{\raggedright\arraybackslash}p{\dimexpr.54\linewidth-3.0pt\relax}@{}}
\toprule
\textbf{Etapa} & \textbf{Resultado confirmado} \\
\midrule
\endfirsthead
\toprule
\textbf{Etapa} & \textbf{Resultado confirmado} \\
\midrule\endhead
1 - Modelos & Sete algoritmos, três dobras nacionais cada; 21 avaliações \\
2 - Seleção & Stacking, SVM linear e Random Forest entre os sete; critério RMSE \\
3 - Curvas iniciais & Três modelos, cinco tamanhos e três dobras; 45 pontos \\
4 - Busca e teste & 60 configurações, 180 avaliações CV, 45 pontos finais e teste de cinco modelos \\
5 - Visualizações & Coleção de 17 gráficos PNG/SVG, 14 tabelas e galeria local \\
6 - Artigo e fonte & Artigo com resultados reais; fonte original reconciliada integralmente \\
7 - Entrega atual & Pitch em roteiro, acompanhamento e artigo em LaTeX; PDFs compilados localmente e pacote Overleaf \\
\bottomrule
\end{longtable}
\endgroup\par\vspace{4pt}
O pedido atual dispensa slides. O pitch é um roteiro oral com divisão nominal das falas, duração planejada de doze minutos, tabela de consulta e perguntas/respostas. Ensaio, apresentação, prazos e submissão permanecem sob responsabilidade da equipe.


\heading{3 Todos os experimentos previstos foram executados?}
Sim, há evidências completas das etapas científicas previstas. Isso não significa que todos os algoritmos foram aplicados ao conjunto de teste. A matriz abaixo explicita o que foi executado e conferido.


\begingroup\singlespacing\setlength{\tabcolsep}{3pt}
\renewcommand{\arraystretch}{1.08}
\begin{longtable}{@{}>{\raggedright\arraybackslash}p{\dimexpr.25\linewidth-4.0pt\relax}>{\raggedright\arraybackslash}p{\dimexpr.35\linewidth-4.0pt\relax}>{\raggedright\arraybackslash}p{\dimexpr.40\linewidth-4.0pt\relax}@{}}
\toprule
\textbf{Tipo de execução} & \textbf{Quantidade e cobertura} & \textbf{Conferência realizada nesta revisão} \\
\midrule
\endfirsthead
\toprule
\textbf{Tipo de execução} & \textbf{Quantidade e cobertura} & \textbf{Conferência realizada nesta revisão} \\
\midrule\endhead
Sete algoritmos solicitados & 7 modelos x 3 dobras = 21 avaliações & Manifestos, hashes, métricas finitas, três dobras e cobertura nacional \\
Referências iniciais & Média, linear, Ridge e árvore; 12 avaliações & Mesmos critérios e partições; 2.153.941 validações por modelo somadas \\
Variantes e hipóteses & 21 avaliações; engenharia, seleção, PCA em dois modelos, ponderação e duas ablações & Resultados por dobra e identidades das fontes \\
Busca de parâmetros & 5 modelos x 12 configurações x 3 dobras = 180 avaliações & Métricas recompostas das previsões persistidas; comparação com GridSearchCV \\
Curvas iniciais & 3 modelos x 5 tamanhos x 3 dobras = 45 pontos & Cobertura, arquivos de pontos e hashes dos ajustes \\
Curvas finais & 45 pontos das receitas congeladas & Identidade dos parâmetros/subconjuntos, integridade e reprodução da CV em 100\% \\
Teste final & 5 modelos; 538.486 participantes iguais por modelo & RMSE, MAE, R\textsuperscript{2} e 260 registros por grupo recompostos das previsões \\
\bottomrule
\end{longtable}
\endgroup\par\vspace{4pt}
Os sete algoritmos são SVM linear, KNN, Bagging, Pasting, Random Forest, Gradient Boosting e Stacking. O conjunto de teste recebeu apenas Stacking, SVM linear, Random Forest, média e linear de referência. KNN, Bagging, Pasting, Boosting, Ridge e árvore isolada não tiveram teste final, conforme o escopo fixado; não se deve ampliar a seleção depois de consultar o teste.


As curvas finais não representam 45 ajustes totalmente novos. Trinta pontos de SVM/Random Forest reaproveitaram configurações idênticas do Checkpoint 3. Doze pontos do Stacking reutilizaram modelos de base com novo combinador treinado sobre previsões fora da dobra interna; três pontos usaram o treino completo das dobras já calculado nas buscas. A origem de cada ponto está registrada. A revisão não treinou novos modelos nem produziu novas previsões do teste.


\heading{4 Resultados e interpretação}
\begingroup\singlespacing\setlength{\tabcolsep}{3pt}
\renewcommand{\arraystretch}{1.08}
\begin{longtable}{@{}>{\raggedright\arraybackslash}p{\dimexpr.31\linewidth-4.5pt\relax}>{\raggedright\arraybackslash}p{\dimexpr.23\linewidth-4.5pt\relax}>{\raggedright\arraybackslash}p{\dimexpr.23\linewidth-4.5pt\relax}>{\raggedright\arraybackslash}p{\dimexpr.23\linewidth-4.5pt\relax}@{}}
\toprule
\textbf{Modelo no teste} & \textbf{RMSE} & \textbf{MAE} & \textbf{R\textsuperscript{2}} \\
\midrule
\endfirsthead
\toprule
\textbf{Modelo no teste} & \textbf{RMSE} & \textbf{MAE} & \textbf{R\textsuperscript{2}} \\
\midrule\endhead
Stacking & 76,190 & 55,325 & 0,2497 \\
Linear de referência & 76,517 & 55,616 & 0,2433 \\
SVM linear & 76,622 & 55,538 & 0,2412 \\
Random Forest & 76,658 & 55,821 & 0,2405 \\
Média do treino & 87,959 & 66,285 & aproximadamente zero \\
\bottomrule
\end{longtable}
\endgroup\par\vspace{4pt}
Fonte: previsões persistidas dos microdados Inep/ENEM 2023. Métricas em 538.486 participantes; RMSE e MAE em pontos. O teste é separado da validação cruzada usada para selecionar receitas.


O Stacking reduz o RMSE da média em 13,38\%. Sua vantagem sobre a linear é pequena e precisa ser ponderada frente ao custo observado. Os resultados do trio são próximos entre CV e teste, mas isso não comprova generalização para outro ENEM. As buscas produziram ganhos mínimos; não foram calculados testes de significância ou intervalos de confiança.


As curvas mostram gaps finais de RMSE de aproximadamente 0,193, 0,007 e 0,400 para Stacking, SVM e Random Forest. Não há sobreajuste acentuado nesses tamanhos, mas R\textsuperscript{2} baixo isoladamente não demonstra subajuste. O menor treino da curva já tem cerca de 143.596 participantes.


Os erros são maiores nos extremos: no Stacking, RMSE de 154,31 abaixo de 400 e 160,74 a partir de 700. Notas baixas são superestimadas e altas são subestimadas em média. Engenharia trouxe ganho pequeno; seleção e PCA pioraram as configurações avaliadas. A ponderação do Ridge piorou o RMSE global para 125,202, embora tenha reduzido erros de extremos na validação. As duas ablações aumentaram o erro, sustentando contribuição preditiva do bloco socioeconômico e da UF, sem interpretação causal.


\heading{5 Testes de software e verificações técnicas}
A suíte existente foi executada nesta revisão: 81 testes aprovados, sem falhas, em 11,44 segundos. Esses testes usam fixtures técnicas de unidade/integração e não constituem experimentos científicos com a população nacional. Não foram criados novos arquivos de testes nem substituídas métricas reais por resultados de fixtures.


\begingroup\singlespacing\setlength{\tabcolsep}{3pt}
\renewcommand{\arraystretch}{1.08}
\begin{longtable}{@{}>{\raggedright\arraybackslash}p{\dimexpr.46\linewidth-3.0pt\relax}>{\raggedright\arraybackslash}p{\dimexpr.54\linewidth-3.0pt\relax}@{}}
\toprule
\textbf{Área da suíte existente} & \textbf{Testes aprovados} \\
\midrule
\endfirsthead
\toprule
\textbf{Área da suíte existente} & \textbf{Testes aprovados} \\
\midrule\endhead
Ingestão & 19 \\
Pré-processamento & 13 \\
Dashboard & 11 \\
Estudo e modelos & 9 \\
Configuração e CLI & 7 \\
Partições & 7 \\
Aquisição e integração & 6 \\
Relações entre atributos & 5 \\
Download oficial & 4 \\
Total & 81 \\
\bottomrule
\end{longtable}
\endgroup\par\vspace{4pt}
Comando executado: \nolinkurl{.venv/bin/python} -m pytest -q, com relatório JUnit em \nolinkurl{resultados/ciencias_natureza/checkpoint-7-2026-11-07/pytest.xml}. A auditoria científica foi executada por \nolinkurl{scripts/auditar_entrega.py} e grava cobertura\_execucoes.csv, validacao\_buscas.json, validacao\_teste.json e auditoria\_experimentos.json no mesmo diretório.


Lint, verificação de formatação e compatibilidade das dependências locais foram aprovados. O ambiente numérico foi mantido; a revisão da entrega não alterou os módulos de treinamento. As fontes LaTeX foram compiladas localmente e os PDFs renderizados para inspeção visual. O pacote foi preparado para importação no Overleaf, com instruções de compilador e arquivo principal.


\heading{6 Limites do que foi validado}
Não foram executados teste temporal em outra edição, validação externa, inferência causal, teste de significância, ensaio oral cronometrado ou apresentação. Também não se declara compilação dentro de uma conta Overleaf: a evidência é a compilação local das fontes fornecidas. A suíte foi executada no Linux deste ambiente; execução Windows depende da automação remota e deve ser consultada separadamente.


Os testes do dashboard verificam componentes existentes e não substituem uma nova inspeção visual interativa de todo o navegador nesta revisão. Há registros históricos de verificações do frontend, mas não foram apresentados como se tivessem sido repetidos agora. A rastreabilidade do experimento usa arquivos reais persistidos e hashes; os hashes locais não são uma assinatura emitida pelo Inep.


\heading{7 Dificuldades, soluções e organização}
A ausência inicial dos originais foi resolvida com a incorporação dos ZIPs fornecidos e comparação integral do CSV. O custo do KNN foi tratado por busca exata compacta, conferida contra distâncias de referência. Interrupções de execução motivaram persistência por dobra e bloco, com retomada condicionada à mesma fonte e configuração. Resultados negativos foram mantidos, como a piora global da ponderação.


A limpeza retirou versões preliminares e pacotes duplicados da área ativa de entrega, preservando um arquivo histórico e o inventário de remoções. Dados originais, banco, partições, modelos ajustados e previsões reais permanecem locais. Arquivos grandes ou individuais não integram o commit; os resultados agregados, código, fontes LaTeX, PDFs e evidências verificáveis são publicados no repositório.


A revisão das alterações precede o commit e o push. A identificação do commit e a confirmação de sincronização com o remoto devem ser consultadas no histórico Git e no relato final da entrega, sem antecipar sucesso antes da operação.


\heading{8 Referências e próximos passos}
Dados oficiais: \url{https://www.gov.br/inep/pt-br/acesso-a-informacao/dados-abertos/microdados/enem}


Repositório: \url{https://github.com/JoeyAlanS/trab2-aprendizado-de-maquina}


Importação no Overleaf: \url{https://docs.overleaf.com/managing-projects-and-files/uploading-a-project}


As referências científicas completas estão no artigo. Para a equipe, faltam ensaiar as falas, conferir o prazo no AVA e realizar a apresentação/submissão. A duração de doze minutos é planejada, não aferida em ensaio.


\end{document}
